\documentclass[11pt]{article}
\usepackage[utf8]{inputenc}
\usepackage[T1]{fontenc}
\usepackage[margin=2.4cm]{geometry}
\usepackage{amsmath,amssymb,amsthm,mathtools}
\usepackage{enumitem}
\usepackage{xcolor}
\usepackage{hyperref}
\hypersetup{colorlinks=true,linkcolor=blue!60!black,citecolor=green!45!black,urlcolor=blue!60!black}

\theoremstyle{plain}
\newtheorem{theorem}{Theorem}[section]
\newtheorem{proposition}[theorem]{Proposition}
\newtheorem{lemma}[theorem]{Lemma}
\newtheorem{corollary}[theorem]{Corollary}
\theoremstyle{definition}
\newtheorem{definition}[theorem]{Definition}
\theoremstyle{remark}
\newtheorem{remark}[theorem]{Remark}

\newcommand{\R}{\mathbb{R}}
\newcommand{\E}{\mathbb{E}}
\newcommand{\Var}{\mathrm{Var}}
\newcommand{\Diff}{\mathrm{Diff}}
\newcommand{\KL}{\mathrm{KL}}
\DeclareMathOperator*{\argmin}{arg\,min}

\title{\textbf{Concentration of Measure under Diffeomorphisms:\\
What Coordinate Choice Can and Cannot Do}\\[0.4em]
\large (corrected and rescoped version)}
\author{Jocelyn Nemb\'e\\[0.2em]
\small L.I.A.G.E, Institut National des Sciences de Gestion, Libreville, Gabon\\
\small Modeling and Calculus Lab, ROOTS-INSIGHTS}
\date{\today}

\begin{document}
\maketitle

\begin{abstract}
We revisit the idea that concentration of measure can be ``improved'' by an
optimal choice of coordinate (diffeomorphism). The transport (covariance) of
concentration inequalities under a diffeomorphism $\phi$ is exact and useful as
an organizing language. However, the stronger claims --- a universal optimal
coordinate, and order-of-magnitude ``strict improvement'' over classical
inequalities --- do not survive scrutiny: the proposed concentration functional
is not scale-normalized, its unconstrained minimization over the diffeomorphism
group is degenerate (infimum $0$), and the advertised improvement factors compare
quantities carried in different units (and, in fact, bounds on \emph{different}
functionals). We give the corrected statements. The genuine residue is twofold:
(i) for a \emph{fixed} target functional the bounded-difference bound is
coordinate-intrinsic, so the only legitimate role of $\phi$ is to \emph{match the
coordinate to the target} (e.g.\ the logarithm for multiplicative data), turning
an inapplicable inequality into an applicable one --- a \emph{qualitative}
(finite-vs-infinite) gain for heavy-tailed/multiplicative data, not a numerical
beat on the same quantity; and (ii) the canonically scale-invariant
concentration object is the Hellinger affinity, multiplicative under products,
whose $-\log$ is additive --- i.e.\ concentration of evidence lives in the
logarithmic twin coordinate, where the Barron--Cover/Kraft bound is the natural
large-deviation estimate. A closing section determines precisely what these
corrected results do, and do not, contribute to orbital/MDL statistical
estimation.
\end{abstract}

\tableofcontents

\section{Introduction and the corrected thesis}

Let $(E,\mathcal A)$ carry a probability measure $\mu$, and let $\phi:E\to F$ be a
diffeomorphism (or measurable bijection) with pushforward $\nu=\phi_*\mu$. The
\emph{transport principle} is the elementary identity
\begin{equation}\label{eq:transport}
\int_E (g\circ\phi)\,d\mu=\int_F g\,d\nu ,\qquad
\mathrm{Law}_\mu\big(g\circ\phi\big)=\mathrm{Law}_\nu(g).
\end{equation}
Thus any concentration statement for the pair $(g,\nu)$ is \emph{the same}
statement for the pair $(g\circ\phi,\mu)$. This is correct and is a convenient
language, but it is an isomorphism: it relabels, it does not create.

The temptation is to treat $\phi$ as a free knob and to ``optimize''
concentration over it. We show this is ill-posed as usually stated and that the
publicized gains are artifacts (Section~\ref{sec:nolunch}). We then isolate the
two things that are genuinely true and useful:
\begin{enumerate}[label=(\roman*)]
\item \textbf{Match the coordinate to the target.} For a \emph{fixed} functional
the bounded-difference / sub-Gaussian constants are intrinsic; the role of $\phi$
is only to express the target where a classical inequality \emph{applies}. For
heavy-tailed or multiplicative data this is a \emph{qualitative} gain (a finite
bound where none existed), not a numerical improvement on the same quantity
(Sections~\ref{sec:multiplicative}--\ref{sec:infogeom}).
\item \textbf{The scale-invariant object is the affinity.} The Hellinger affinity
is multiplicative under products, so $-\log$ of it is additive: evidence
concentrates in the logarithmic (twin) coordinate, where the Barron--Cover/Kraft
bound is the natural large-deviation estimate (Section~\ref{sec:link}).
\end{enumerate}

\section{The covariance dictionary (correct)}
\label{sec:dictionary}

\begin{definition}[Pushforward, pullback]\label{def:pp}
For $\phi:E\to F$ and $\mu$ on $E$, the pushforward is
$(\phi_*\mu)(B)=\mu(\phi^{-1}(B))$; for $g:F\to\R$ the pullback is
$\phi^*g=g\circ\phi$; for a metric $d_F$ the pullback metric is
$(\phi^*d_F)(x,y)=d_F(\phi(x),\phi(y))$.
\end{definition}

\begin{lemma}[Exact covariance of concentration]\label{lem:cov}
Let $\nu=\phi_*\mu$. Then for every functional $g$,
$\mathrm{Law}_\mu(g\circ\phi)=\mathrm{Law}_\nu(g)$. Consequently:
\begin{enumerate}[label=(\alph*)]
\item $g\circ\phi$ is sub-Gaussian under $\mu$ with parameter $\sigma^2$ iff $g$
is sub-Gaussian under $\nu$ with the same parameter $\sigma^2$;
\item $\nu$ satisfies a log-Sobolev (resp.\ Poincar\'e) inequality with constant
$C$ iff $\mu$ satisfies the pulled-back inequality with the \emph{same} constant
$C$;
\item the isoperimetric profile is covariant: $I_\mu^{\phi^*d}=I_\nu^{d}$.
\end{enumerate}
\end{lemma}

\begin{proof}
Immediate from \eqref{eq:transport}: entropy, variance, Lipschitz seminorms and
boundary measures are computed through integrals that are invariant under the
change of variables $x\mapsto\phi(x)$.
\end{proof}

\begin{remark}[What Lemma~\ref{lem:cov} does \emph{not} say]
It transports a \emph{single} inequality between a measure and its pushforward.
It does \emph{not} unify the distinct classical inequalities (Hoeffding,
Bernstein, Talagrand, Gaussian isoperimetry): these rest on different hypotheses
(bounded range, variance control, convexity, Gaussianity) and are not coordinate
images of one another. The earlier claim that they are ``manifestations of a
single principle'' over-reaches.
\end{remark}

\section{No free improvement: the central correction}
\label{sec:nolunch}

Fix the sub-Gaussian functional of a measure,
\begin{equation}\label{eq:F}
F[\mu,\phi]:=\inf\{\sigma^2:\ \phi_*\mu\ \text{is sub-Gaussian with parameter}\ \sigma^2\}.
\end{equation}

\begin{proposition}[The unconstrained optimum is degenerate]\label{prop:degenerate}
For any non-degenerate $\mu$ on $\R^n$,
\[
\inf_{\phi\in\Diff(\R^n)} F[\mu,\phi]=0,
\]
and the infimum is not attained by a diffeomorphism: it is approached by
compressive maps $\phi=\varepsilon\,\psi$ ($\psi$ bounded, $\varepsilon\to0$), or by any
$\phi$ with $\phi_*\mu$ supported in a vanishing ball. Hence
``$\phi^\star=\argmin_\phi\sup_{\mu}F[\mu,\phi]$'' is \emph{ill-posed} without a
normalization of $\phi$.
\end{proposition}

\begin{proof}
$F[\varepsilon\psi\cdot\,;\mu]=\varepsilon^2 F[\psi_*\mu]\to0$. A deviation of
$\phi(X)$ by $t$ corresponds to a deviation of $X$ by $\asymp t/\varepsilon\to\infty$;
the small ``concentration constant'' is bought by inflating the scale of the
quantity being controlled.
\end{proof}

The correct way to make the question meaningful is to \emph{fix what is being
bounded}.

\begin{proposition}[Bounded-difference bounds are coordinate-intrinsic]\label{prop:intrinsic}
Let $h:E^n\to\R$ be a fixed functional with bounded differences $c_i$: for all
$x,x_i'$,
$|h(\dots,x_i,\dots)-h(\dots,x_i',\dots)|\le c_i$ whenever $x_i,x_i'$ range over
the support. Then $P(|h-\E h|\ge t)\le 2\exp(-2t^2/\sum_i c_i^2)$, and the
constants $c_i$ depend only on $h$ and the supports, \emph{not} on any coordinate
$\phi$ used to describe the data. In particular no diffeomorphism improves the
McDiarmid bound for a fixed $h$.
\end{proposition}

\begin{proof}
McDiarmid's inequality applied to $h$. Reparametrizing the data by $\phi$ replaces
$h$ by $h\circ\phi^{-1}$ acting on $\phi(X_i)$; the increments are unchanged
because $h\circ\phi^{-1}(\phi(x))=h(x)$, so $c_i$ is the same number computed in
either chart.
\end{proof}

\begin{remark}[Consequence for ``optimal coordinates'']
Propositions~\ref{prop:degenerate}--\ref{prop:intrinsic} together say: the only
legitimate effect of $\phi$ is to change \emph{which functional} one bounds (or
to make a previously inapplicable inequality applicable). There is no universal
coordinate that tightens concentration for a fixed target; the well-posed problem
is to \emph{match} $\phi$ to the target functional.
\end{remark}

\section{Multiplicative data: the matched coordinate (corrected)}
\label{sec:multiplicative}

Let $X_1,\dots,X_n$ be independent with $X_i\in[a_i,b_i]\subset(0,\infty)$, and
let $P_n=\prod_i X_i$. The natural (scale-invariant) deviation of a product is
\emph{multiplicative}, i.e.\ a deviation of $\log P_n$.

\begin{theorem}[Concentration of products, corrected]\label{thm:product}
\[
P\big(|\log P_n-\E\log P_n|\ge t\big)\ \le\ 2\exp\!\Big(-\frac{2t^2}{\sum_{i=1}^n\log^2(b_i/a_i)}\Big).
\]
Equivalently, the multiplicative deviation $P_n/\exp(\E\log P_n)$ is controlled
at the dimensionless scale $\sum_i\log^2(b_i/a_i)$.
\end{theorem}

\begin{proof}
$\log P_n=\sum_i\log X_i$ with $\log X_i\in[\log a_i,\log b_i]$ of range
$\log(b_i/a_i)$; apply Hoeffding. By Proposition~\ref{prop:intrinsic} the same
constant is obtained by McDiarmid on the functional $\log P_n$ in \emph{any}
chart: the logarithm is merely the chart in which $\log P_n$ is a sum of bounded
terms.
\end{proof}

\begin{remark}[Why there is no ``$\Gamma\approx 21000$'' free lunch]
The previously claimed improvement factor $\Gamma=(b-a)^2/(\log(b/a))^2$ compares
two \emph{different} objects: the Hoeffding constant for the \emph{sum}
$S_n=\sum_iX_i$ (in units of $X^2$) against the Hoeffding constant for the
\emph{product} $P_n$ (dimensionless). The numerator is scale-dependent (under
$X\mapsto\lambda X$ it is multiplied by $\lambda^2$), the denominator is not, so
$\Gamma$ can be made arbitrary by rescaling and is not a function of $b/a$ alone.
For one \emph{fixed} functional there is no such factor
(Proposition~\ref{prop:intrinsic}).
\end{remark}

\begin{theorem}[Qualitative gain for heavy tails]\label{thm:heavy}
Suppose $X>0$ is heavy-tailed: it admits no finite sub-Gaussian (or
Bernstein-type) parameter in the identity chart, while $\log X$ is
bounded/sub-Gaussian. Then no Hoeffding/Bernstein bound exists for $X$ directly,
whereas Theorem~\ref{thm:product} and Lemma~\ref{lem:cov} yield finite
concentration for the multiplicative functionals of $X$ (products, geometric
means, quantiles). The improvement is \emph{qualitative} (finite versus infinite),
not a numerical ratio on a common quantity.
\end{theorem}

\begin{proof}
Direct from Lemma~\ref{lem:cov}: sub-Gaussianity of $\log X$ transports to all
functionals expressible through $\log X$; no statement about $X$ in the identity
chart is implied or claimed.
\end{proof}

\begin{theorem}[Worst case over a bounded class, corrected]\label{thm:minimax}
Over all $\mu$ supported in $[a,b]$, the sub-Gaussian functional is maximized,
not minimized, at the two-point (Rademacher) law on $\{a,b\}$:
\[
\sup_{\mathrm{supp}(\mu)\subseteq[a,b]} F[\mu,\mathrm{id}]=\frac{(b-a)^2}{4},
\qquad
\inf_{\mathrm{supp}(\mu)\subseteq[a,b]} F[\mu,\mathrm{id}]=0\ (\text{Dirac}).
\]
Hence the \emph{class-uniform} Hoeffding constant in the identity chart is
$(b-a)^2/4$ (a $\sup$), and in the logarithmic chart $(\log(b/a))^2/4$; these
are worst-case constants for \emph{different} functionals/units and are not
directly comparable.
\end{theorem}

\begin{proof}
For a distribution on $[a,b]$, $\Var\le (b-a)^2/4$ with equality at the endpoints
two-point law (Popoviciu/Bhatia--Davis), and $\Var\ge0$ with equality at a Dirac.
The earlier statement labelled this extremum an ``$\inf$''; it is a $\sup$.
\end{proof}

\section{Information-geometric coordinate (corrected)}
\label{sec:infogeom}

\begin{proposition}[Expectation coordinate as the matched chart]\label{prop:expfam}
Let $P_\theta(dx)=\exp(\langle\theta,T(x)\rangle-A(\theta))\,\mu_0(dx)$ be a
regular exponential family with log-partition $A$. The expectation coordinate
$\eta=\nabla A(\theta)=\E_\theta[T]$ is the canonical chart for the sufficient
statistic: the Fisher information $\nabla^2 A(\theta)=\Var_\theta(T)$ transforms
into $\nabla^2 A^\ast(\eta)$ with $A^\ast$ the Legendre dual, and the family is
\emph{dually flat} (the $e$- and $m$-connections are flat).
\end{proposition}

\begin{proof}
Standard convex-duality computation; $\eta=\nabla A(\theta)$, $\theta=\nabla
A^\ast(\eta)$, $\nabla^2A^\ast=(\nabla^2A)^{-1}$.
\end{proof}

\begin{remark}[Correction to ``Fisher--Rao becomes Euclidean'']
Dual flatness is \emph{not} the same as making the Fisher--Rao (Levi-Civita)
metric Euclidean. For $\dim\ge 2$ the Fisher--Rao metric of an exponential family
generally has nonzero Riemannian curvature, so \emph{no} chart renders it
Euclidean; only the dual-affine structure flattens. The matched chart $\eta$ is
canonical in the dually-flat sense, which is the correct statement.
\end{remark}

\section{The scale-invariant object: affinity concentration and the link to estimation}
\label{sec:link}

The well-posed, scale-invariant concentration object is not $F[\mu,\phi]$ but the
\emph{Hellinger affinity}
\[
\rho(P,Q)=\int\sqrt{dP\,dQ}\in(0,1],\qquad h^2(P,Q)=1-\rho(P,Q).
\]

\begin{proposition}[Evidence concentrates in the logarithmic twin coordinate]\label{prop:affinity}
The affinity is multiplicative under products, $\rho(P^{\otimes n},Q^{\otimes
n})=\rho(P,Q)^n$, so $-\log\rho$ is additive. For a countable model family
$\{f_m\}$ with prefix code lengths satisfying the half-Kraft condition
$\sum_m 2^{-L(m)/2}\le 1$, the minimum-description-length selector $\widehat f$
obeys the partition-function (large-deviation) bound
\begin{equation}\label{eq:bc}
\E_{P^\star}\!\Big[2^{-L(\widehat m)/2}\textstyle\prod_{i=1}^n\sqrt{\widehat f(X_i)/f^\star(X_i)}\Big]
\ \le\ \sum_m 2^{-L(m)/2}\rho(f^\star,f_m)^n\ \le\ 1,
\end{equation}
whence $\E[h^2(P^\star,\widehat f)]\le\inf_m\{\KL(P^\star\|f_m)+(\log2)L(m)/n\}$.
\end{proposition}

\begin{proof}
Multiplicativity of $\rho$ is Fubini; \eqref{eq:bc} is the Barron--Cover
argument: $\widehat m$ maximizes $2^{-L(m)}\prod_i f_m(X_i)$, take square roots
and bound the max by the sum; then take $\E_{P^\star}$. The resolvability bound
follows by the standard passage (Jensen, $h^2\le-\log\rho\le\KL$).
\end{proof}

This is the correct form of ``concentration in a twin space'': the relevant
non-additive structure is the \emph{multiplicative} space of affinities/likelihood
ratios, whose additive (logarithmic) chart is canonical and \emph{scale
invariant} --- precisely the normalization that $F[\mu,\phi]$ lacks.

\begin{remark}[Scope: near-minimax, not exact]
The resolvability bound of Proposition~\ref{prop:affinity} is the genuine
contribution to orbital/MDL estimation, but it must be read precisely. On a
nonparametric sieve the code length is $L(m)\asymp D_m\log n$, so the bound yields
the rate $(\log n/n)^{2s/(2s+d)}$ --- minimax \emph{up to a logarithmic factor},
the per-coefficient quantization cost of the description-length code. The
affinity/MDL machinery is thus \emph{near-minimax}; the \emph{exact} rate
$n^{-2s/(2s+d)}$ is attained not by this code but by a projection /
penalized-projection estimator paying the metric complexity $D_m/n$. What these
corrected results contribute is the concentration step (the affinity bound), not
the removal of the log, which is a separate, estimator-level refinement.
\end{remark}

\section{Honest scope: what is, and is not, a contribution}
\label{sec:scope}

\noindent\textbf{Retained / correct.}
\begin{itemize}[leftmargin=1.4em]
\item The covariance dictionary (Lemma~\ref{lem:cov}): a clean transport of
sub-Gaussianity, log-Sobolev, Poincar\'e, isoperimetry between a measure and its
pushforward.
\item The matched-coordinate principle for multiplicative/heavy-tailed data
(Theorems~\ref{thm:product}--\ref{thm:heavy}): a \emph{qualitative} extension of
applicability.
\item The expectation/dually-flat chart for exponential families
(Proposition~\ref{prop:expfam}).
\item The affinity formulation (Proposition~\ref{prop:affinity}) as the
scale-invariant concentration object.
\end{itemize}

\noindent\textbf{Withdrawn / corrected.}
\begin{itemize}[leftmargin=1.4em]
\item No universal optimal coordinate: $\inf_\phi F[\mu,\phi]=0$
(Proposition~\ref{prop:degenerate}).
\item No order-of-magnitude ``strict improvement'': the factor
$(b-a)^2/(\log(b/a))^2$ mixes units and functionals
(Proposition~\ref{prop:intrinsic}, Theorem~\ref{thm:product}).
\item The internal inconsistency between a threshold $b/a=e^2$ and a uniformly
$\ge1$ improvement factor is removed (there is no such factor for a fixed
functional).
\item Minimax $\inf$ corrected to $\sup$ (Theorem~\ref{thm:minimax}).
\item No unification of Hoeffding/Bernstein/Talagrand/Gaussian isoperimetry; no
chart making Fisher--Rao Euclidean.
\end{itemize}

\begin{thebibliography}{9}
\bibitem{boucheron} S.~Boucheron, G.~Lugosi, P.~Massart, \emph{Concentration
Inequalities}, Oxford Univ.\ Press, 2013.
\bibitem{barroncover} A.~Barron, T.~Cover, ``Minimum complexity density
estimation,'' \emph{IEEE Trans.\ Inform.\ Theory} 37 (1991) 1034--1054.
\bibitem{amari} S.~Amari, \emph{Information Geometry and Its Applications},
Springer, 2016.
\bibitem{villani} C.~Villani, \emph{Optimal Transport: Old and New}, Springer,
2009.
\end{thebibliography}

\end{document}
